% !TeX root = ../../Book2.tex
% 纲要标题：## 第9章　外代数与行列式
% 纲要位置：计划纲要.md，第 1066 行。

\chapter{外代数与行列式}
\label{b2:ch:09}
\BookChapterNumber{b2:ch:09}

\begin{chapterabstract}
外幂把交替多线性映射化为线性映射，外代数则把不同次数的交替运算统一起来。本章先区分交替性与反对称性，在任意交换环上构造外幂并证明有限自由模的外幂基定理；随后从最高外幂定义行列式，推导乘法性与 Cauchy–Binet 公式。伴随矩阵、Cayley–Hamilton 定理、四维空间中的 Plücker 关系以及 Koszul 微分，分别展示外代数在线性算子、子空间坐标和正合性问题中的作用。
\end{chapterabstract}

\begin{learninggoals}
\begin{itemize}
\item 分清交替性与反对称性，说明特征二或目标模中的二挠为何影响它们的等价性。
\item 用泛性质构造外幂和楔积，证明外幂基、分次交换律及直和分解，不借助除以阶乘的公式。
\item 从行列式线推导行列式乘法性与子式公式，并在交换环上证明伴随矩阵恒等式和 Cayley–Hamilton 定理。
\item 在四维向量空间中用 Plücker 坐标判定二向量是否可分解，写出收缩算子和 Koszul 微分，区分复合为零与正合。
\end{itemize}
\end{learninggoals}

在二维自由模的基 \(e_1,e_2\) 下，令 \(v=ae_1+ce_2\)、\(w=be_1+de_2\)。我们希望构造一个双线性运算，记作 \(v\wedge w\)，要求两个相同向量的积为零。这样的交替运算在交换输入时变号，按双线性展开后只剩
\[
v\wedge w=(ad-bc)e_1\wedge e_2.
\]
系数 \(ad-bc\) 正是二阶行列式。此时楔积还只是待构造运算的记号；这个计算说明交替双线性规则会产生面积式的系数。下面用张量积的商实现这种运算，再证明最高外幂上的线性作用恰由这个标量描述。

\ProseParagraphBreak

要把这种计算推广到任意交换环，不能预先把“交换时变号”当作“同一向量的楔积为零”的同义说法：在特征二它们并不等价。从交替关系构造外幂后，基和符号规则都须由这些关系推出。这样得到的运算还可用于多行多列的子式，以及线性算子所满足的多项式关系。

\ProseParagraphBreak

本章需要第六章的张量积泛性质与第八章的多重张量、张量代数。主要构造允许任意交换环；只在把楔积与线性无关、子空间相联系时改用域。行列式的坐标公式由外幂的交替求和得到。

\ProseParagraphBreak
松村英之的《Commutative Ring Theory》\cite[Appendix C, (1)--(6), pp. 283--286]{Matsumura1989CommutativeRingTheory}讨论一般交换环上的外幂、直和与收缩，适合与本章的泛性质证明对照。Roman 的《Advanced Linear Algebra》\cite[Chapter 14, The Determinant, pp. 403--405]{Roman2008AdvancedLinearAlgebra}从多线性形式解释行列式；交替与反对称的定义在不同特征下须加以区别。

% !TeX root = ../../Book2.tex
% 纲要标题：### 9.1 交替多线性映射
% 纲要位置：计划纲要.md，第 1068 行。

\section{交替多线性映射}
\label{b2:sec:0901}

面积在交换两条边时改变符号，在两条边重合时变为零。这两种性质在熟悉的实数情形中相互联系，但在一般系数环上并不等价。外幂所表达的是“有两个输入相同时为零”的交替性。本章始终设 \(R\) 为含单位元的交换环，所有张量积和多线性映射均在 \(R\) 上；需要域或额外特征条件时另行说明。

% 纲要要点（供目录核对）：
%
% * 交替性与反对称性的区别；双线性矩阵的对角线条件
% * 特征二与目标模中的二挠现象
% * 外幂的泛性质；先用张量积表达多线性，再按重复输入取商
% * 外幂不保持一般单射的理想反例，完整证明放在正文
%

\subsection{交替性与反对称性}
\label{b2:subsec:0901:01}

\begin{definition}\label{b2:def:alternating-multilinear}
设 \(R\) 为交换环，\(M,L\) 为 \(R\) 模，\(p\ge1\)。对逐变量 \(R\) 线性的映射 \(b:M^p\to L\)，若任意两个输入相同便有
\[
b(m_1,\ldots,m_p)=0,
\]
就称 \(b\) 为\term[jiaoti-duoxianxing-yingshe]{交替多线性映射}[alternating multilinear map]。若对任意两个位置，交换输入后 \(b\) 的值都变为原值的负元，就称 \(b\) 为\term[fanduichen-duoxianxing-yingshe]{反对称多线性映射}[skew-symmetric multilinear map]。\(p=1\) 时，两种条件均为空条件，此时这两类映射就是线性映射。
\end{definition}

\begin{proposition}\label{b2:prop:alternating-implies-skew}
设 \(R\) 为交换环，\(M,L\) 为 \(R\) 模，\(p\ge1\)。从 \(M^p\) 到 \(L\) 的交替多线性映射必反对称。反过来，若目标模 \(L\) 中 \(2\ell=0\) 蕴含 \(\ell=0\)，则从 \(M^p\) 到 \(L\) 的反对称多线性映射也交替。特别地，当 \(2\) 在 \(R\) 中可逆时，两种条件等价。
\end{proposition}
\begin{proof}
\(p=1\) 时两个条件均为空条件。以下设 \(p\ge2\)。
固定除第 \(i,j\) 个位置以外的所有输入，把这两个位置都取为 \(x+y\)。交替性使函数值为零；按这两个变量展开，两个含 \((x,x)\)、\((y,y)\) 的项分别为零，故剩下的两项满足
\[
b(\ldots,x,\ldots,y,\ldots)
+b(\ldots,y,\ldots,x,\ldots)=0.
\]
这就是反对称性。

若反对称映射在两个位置取相同的 \(x\)，交换它们不改变输入，却使值变号，所以这个值 \(\ell\) 满足 \(2\ell=0\)。所给假设使 \(\ell=0\)，即交替。当 \(2\) 可逆时，乘以 \(2^{-1}\) 就能从 \(2\ell=0\) 推出 \(\ell=0\)。
\end{proof}

将置换写成对换的复合，交替性还给出
\begin{equation}\label{b2:eq:alternating-permutation}
b(m_{\sigma(1)},\ldots,m_{\sigma(p)})
=\operatorname{sgn}(\sigma)b(m_1,\ldots,m_p)
\qquad(\sigma\in S_p).
\end{equation}
所用的符号同态及其在换位上取值为 \(-1\) 的性质，已由第一卷定理3.2.3证明。符号 \(\operatorname{sgn}(\sigma)\) 只由置换决定，因此表达不依赖所选的对换分解；公式在 \(p!\) 不可逆的环上也成立。

\ProseParagraphBreak
例如，任意交换环上的
\[
b\bigl((x_1,x_2),(y_1,y_2)\bigr)=x_1y_2-x_2y_1
\]
是交替双线性映射。其交替性应直接由 \(x_1x_2-x_2x_1=0\) 核对；不能先用“交换变号”，再在可能存在二挠的环上消去一个 \(2\)。

\subsection{特征二与二挠现象}
\label{b2:subsec:0901:02}

在域 \(k=\mathbb F_2\) 上，双线性映射 \(b:k\times k\to k\)、\(b(x,y)=xy\) 满足
\[
b(y,x)=b(x,y)=-b(x,y),
\]
所以反对称，却有 \(b(1,1)=1\)，因而不交替。在特征二中，“交换变号”变成“交换不变”，完全没有强迫对角线取值为零。

\ProseParagraphBreak
这一区别不只由环的特征决定。取 \(R=\mathbb Z\)、\(M=\mathbb Z\)、\(L=\mathbb Z/2\mathbb Z\)，公式 \(b(a,c)=\overline{ac}\) 同样反对称而不交替。环本身特征零，目标模却有二挠。因此，使用\cref{b2:prop:alternating-implies-skew}从反对称性推出交替性时，要检查的是实际取值所在的目标模能否消去 \(2\)，而非仅仅检查底环的特征。

\ProseParagraphBreak
对于交换环 \(R\) 上基为 \(e_1,\ldots,e_n\) 的有限自由模 \(M\)，一个双线性形式 \(b:M\times M\rightarrow R\) 由矩阵 \(A=(a_{ij})\)、\(a_{ij}=b(e_i,e_j)\) 确定。反对称性等价于 \(A^{\mathsf t}=-A\)，而交替性还要求对角线为零，即
\[
a_{ii}=0,\qquad a_{ij}=-a_{ji}.
\]
必要性由交替性及\cref{b2:prop:alternating-implies-skew}得到。反过来，对 \(x=\sum_i x_ie_i\)，展开
\[
b(x,x)=\sum_i a_{ii}x_i^2+
\sum_{i<j}(a_{ij}+a_{ji})x_ix_j=0.
\]
所以条件充分。特征二中，这种矩阵是对称且对角线为零的矩阵；一般对称矩阵并不足够。后续讨论交替双线性型时，仍以对角线取值为零作为交替性的定义。

\subsection{外幂的泛性质}
\label{b2:subsec:0901:03}

前面的区别决定了外幂应当施加哪些关系。如果在二重张量积中只把 \(x\otimes y+y\otimes x\) 置为零，特征二时这只强制交换不变，并不强制 \(x\otimes x=0\)。例如 \(M=\mathbb F_2 e\) 时，所有这样的反对称关系本来就是零，非零的 \(e\otimes e\) 会被保留下来。因此，要表达交替性，必须直接把重复输入所对应的张量置为零。

\begin{definition}\label{b2:def:exterior-power}
设 \(R\) 为交换环，\(M\) 为 \(R\) 模。对 \(p\ge2\)，在 \(p\) 重张量积 \(M^{\otimes p}\) 中，令 \(J_p\) 为所有具有两个相同因子的纯张量生成的子模。将商模
\[
\bigwedge^p M\coloneqq M^{\otimes p}/J_p,
\]
称为 \(M\) 的 \(p\) 次\term[waimi]{外幂}[exterior power]，并把 \(m_1\otimes\cdots\otimes m_p\) 的像记为 \(m_1\wedge\cdots\wedge m_p\)。约定 \(\bigwedge^0M=R\)、\(\bigwedge^1M=M\)，空楔积为 \(1\in R\)。
\end{definition}
\symidx{\(\bigwedge^p M\)：模的第 \(p\) 次外幂}
\symidx{\(m_1\wedge\cdots\wedge m_p\)：外幂中的纯楔积}

多重张量积已经把 \(p\) 个输入的多线性关系转化为模中的线性关系；再模掉 \(J_p\)，就是在这个模中加入“重复输入为零”的关系。这与第三章按关系取模的商、以及第八章通过张量商施加代数关系的构造相同。所得映射 \((m_1,\ldots,m_p)\mapsto m_1\wedge\cdots\wedge m_p\) 多线性且交替，所有这些纯楔积生成外幂。

\begin{proposition}[外幂的泛性质]\label{b2:prop:exterior-alternating-universal}
设 \(R\) 为交换环，\(M,L\) 为 \(R\) 模，\(p\ge1\)。对任意交替 \(p\) 线性映射 \(b:M^p\to L\)，存在唯一 \(R\) 线性映射
\[
\widetilde b:\bigwedge^pM\longrightarrow L,
\qquad \widetilde b(m_1\wedge\cdots\wedge m_p)=b(m_1,\ldots,m_p).
\]
反过来，每个这样的线性映射都给出交替 \(p\) 线性映射。具有这一性质的对象连同其指定映射，在唯一的相容同构意义下确定。
\end{proposition}
\begin{proof}
\(p=1\) 时，\(\bigwedge^1M=M\)，结论就是线性映射本身。以下设 \(p\ge2\)。由\cref{b2:thm:multiple-tensor-universal}，多线性映射 \(b\) 唯一延伸为 \(M^{\otimes p}\to L\)。交替性恰好说明 \(J_p\) 的每个生成元都被送到零，所以它唯一地通过商模分解。反过来，商映射本身交替，与线性映射复合仍然交替。唯一性来自纯楔积的生成性。

若另一个对象也具有同样的性质，两次应用泛性质得到双向线性映射；它们的复合在每个指定纯楔积上为恒等，故处处为恒等，证明唯一同构。\(p=0\) 时把零变量的值看成指定的 \(\ell\in L\)，它唯一对应 \(R\to L\)、\(r\mapsto r\ell\)。
\end{proof}

因此，要定义一个从 \(\bigwedge^pM\) 到 \(L\) 的线性映射，可以先给出 \(M^p\) 上的公式，检查它对各变量线性且在重复输入时为零，再由泛性质取得唯一的线性映射。后面读取外幂坐标、定义楔积及收缩算子时，使用的都是这一方法。

\ProseParagraphBreak
设 \(f:M\to N\) 线性。交替公式
\[
(m_1,\ldots,m_p)\longmapsto
f(m_1)\wedge\cdots\wedge f(m_p)
\]
诱导 \(\bigwedge^p f:\bigwedge^pM\to\bigwedge^pN\)。在纯楔积上计算得到
\begin{equation}\label{b2:eq:exterior-functoriality}
\bigwedge^p(gf)=\left(\bigwedge^p g\right)\left(\bigwedge^p f\right),
\qquad \bigwedge^p\operatorname{id}_M=\operatorname{id}_{\bigwedge^pM}.
\end{equation}
于是每个外幂都是一个函子；\(p=0\) 时规定所有 \(\bigwedge^0 f\) 都为 \(\operatorname{id}_R\)。与张量积相似，纯楔积的表达也不唯一，外幂的一般元素是有限和，未必能写成一个纯楔积。

\ProseParagraphBreak
函子化保证每个线性映射都诱导外幂上的映射，却不保证它保留单射。由张量积取商得到外幂，并没有消除第七章讨论过的这种障碍。

\begin{proposition}\label{b2:prop:exterior-injection-failure}
设 \(k\) 为域，\(R=k[x,y]\)，\(I=(x,y)\) 为 \(R\) 的理想，并把 \(I\) 视为 \(R\) 模。包含映射 \(i:I\hookrightarrow R\) 是单射，但诱导映射
\[
\bigwedge^2i:\bigwedge^2_R I\longrightarrow\bigwedge^2_R R
\]
不是单射：\(x\wedge y\ne0\) 于 \(\bigwedge^2_R I\)，而其像为零。
\end{proposition}
\begin{proof}
为了检测 \(x\wedge y\)，只需构造一个在 \((x,y)\) 处不为零的交替双线性映射。把 \(k\) 视为经 \(R\rightarrow R/I\cong k\)、\(r\mapsto r(0,0)\) 取得标量作用的 \(R\) 模。任意 \(f,g\in I\) 唯一地写成
\[
f=ax+by+f_{\ge2},\qquad g=cx+dy+g_{\ge2},
\qquad a,b,c,d\in k,
\]
其中余下的项全次数至少为二。定义 \(\beta(f,g)=ad-bc\)。读取一次项系数是可加的，所以 \(\beta\) 对两个变量均可加。对任意 \(r\in R\)，\(rf\) 的一次项恰为 \(r(0,0)(ax+by)\)，因此
\[
\beta(rf,g)=r(0,0)\beta(f,g),
\qquad \beta(f,rg)=r(0,0)\beta(f,g).
\]
这就是以 \(k\cong R/I\) 为目标的 \(R\) 双线性。又 \(\beta(f,f)=ab-ba=0\)，所以 \(\beta\) 交替。由\cref{b2:prop:exterior-alternating-universal}，它诱导线性映射 \(\widetilde\beta:\bigwedge^2_RI\rightarrow k\)，满足 \(\widetilde\beta(x\wedge y)=\beta(x,y)=1\)。于是 \(x\wedge y\ne0\)。

另一方面，\(R\) 作为自身的模由 \(1\) 生成；任意纯楔积 \(r\wedge s=rs(1\wedge1)=0\)。纯楔积生成二次外幂，所以 \(\bigwedge^2_R R=0\)。故 \(\bigwedge^2i\) 将上述非零元素送到零，不是单射。
\end{proof}

% !TeX root = ../../Book2.tex
% 纲要标题：### 9.2 外代数
% 纲要位置：计划纲要.md，第 1074 行。

\section{外代数}
\label{b2:sec:0902}

第八章的对称代数把生成元交换前后的差规定为零；外代数则把每个一次元素的平方规定为零，由此得到反对称关系。两者都是张量代数的商，却分别对应对称与交替的多线性运算。即使在特征二中，外代数的平方为零关系也不会消失。

\ProseParagraphBreak
不同次数的外幂不是彼此孤立的模。把输入首尾连接，可以在它们之间定义楔积，组成外代数。有基时，其计算归结为把指标排列成递增次序；没有基时，构造仍由交替映射的泛性质保证。本节把这两种描述联系起来，并核对它们在任意特征下都成立。

% 纲要要点（供目录核对）：
%
% * 基、维数与楔积；用交替坐标函数检测外幂基的系数
% * 分次交换律与奇数次平方为零；代数泛性质要求所有一次元素平方为零
% * 直和的外代数；分次张量积的符号与按次数分解
% * 域上纯楔积非零与线性无关的判据；一般环上非零与成基的区别
%

\subsection{基、维数与楔积}
\label{b2:subsec:0902:01}

设 \(\xi\in\bigwedge^pM\)、\(\eta\in\bigwedge^qM\)。希望用下式定义乘法：
\[
(m_1\wedge\cdots\wedge m_p)\wedge(n_1\wedge\cdots\wedge n_q)
=m_1\wedge\cdots\wedge m_p\wedge n_1\wedge\cdots\wedge n_q.
\]
纯楔积并非自由符号，因而需要先证明这个公式与它们之间的关系相容。

\begin{proposition}\label{b2:prop:wedge-product}
设 \(R\) 为交换环，\(M\) 为 \(R\) 模，\(p,q\ge0\)。规定两个纯楔积的积为
\[
(m_1\wedge\cdots\wedge m_p)\wedge(n_1\wedge\cdots\wedge n_q)
=m_1\wedge\cdots\wedge m_p\wedge n_1\wedge\cdots\wedge n_q.
\]
这一公式唯一延伸为 \(R\) 双线性映射
\[
\bigwedge^pM\times\bigwedge^qM\longrightarrow\bigwedge^{p+q}M,
\qquad(\xi,\eta)\longmapsto\xi\wedge\eta.
\]
这些映射满足结合律，\(R=\bigwedge^0M\) 中的 \(1\) 为单位元。因此
\[
\bigwedge M\coloneqq\bigoplus_{p\ge0}\bigwedge^pM
\]
成为分次 \(R\) 代数，称为 \(M\) 的\term[waidaishu]{外代数}[exterior algebra]；其乘法称为\term[xieji]{楔积}[wedge product]。
\end{proposition}
\begin{proof}
在 \(p=2,q=1\) 时，固定 \(n\in M\)，公式 \((m_1,m_2)\mapsto m_1\wedge m_2\wedge n\) 是交替双线性的。由\cref{b2:prop:exterior-alternating-universal}，它唯一地给出线性映射 \(L_n:\bigwedge^2M\rightarrow\bigwedge^3M\)，满足
\[
L_n(m_1\wedge m_2)=m_1\wedge m_2\wedge n.
\]
对 \(r,s\in R\) 及 \(n,n'\in M\)，\(L_{rn+sn'}=r L_n+s L_{n'}\) 在纯楔积上由多线性成立，故在整个 \(\bigwedge^2M\) 上成立。因此 \((\xi,n)\mapsto L_n(\xi)\) 双线性，且不依赖 \(\xi\) 的纯楔积表达。这就是把前两个变量先合成外幂，再对第三个变量使用一次外幂 \(\bigwedge^1M=M\) 的泛性质。

一般地，先设 \(p,q\ge1\)。固定 \(n_1,\ldots,n_q\)，右端关于 \(m_1,\ldots,m_p\) 交替且多线性，所以对 \(\xi\) 唯一线性延伸。所得公式关于各 \(n_j\) 仍多线性且交替：在纯 \(\xi\) 上已有这些等式，再由线性延伸到所有 \(\xi\)。第二次应用泛性质，即得关于 \(\eta\) 的唯一线性延伸。

若任一次数为零，则定义 \(r\wedge\xi=\xi\wedge r=r\cdot\xi\)，其中 \(r\in R=\bigwedge^0M\)。结合律在三个正次纯楔积上只是同一串因子的首尾连接，三线性与生成性使其在任意三个元素上成立；含零次因子时则由标量作用的结合律成立。\(1\) 因而为单位。直和中的元素只有有限多个非零次数分量，乘法扩张后的求和仍有限，故给出整个直和上的代数结构。
\end{proof}
\symidx{\(\bigwedge M\)：模 \(M\) 的外代数}

\begin{theorem}[有限自由模的外幂基]\label{b2:thm:exterior-power-basis}
设 \(R\) 为交换环，\(M\) 为有基 \(e_1,\ldots,e_n\) 的有限自由 \(R\) 模。对 \(0\le p\le n\)，全部递增指标楔积
\[
e_I\coloneqq e_{i_1}\wedge\cdots\wedge e_{i_p},
\qquad I=\{i_1<\cdots<i_p\}\subseteq\{1,\ldots,n\},
\]
组成 \(\bigwedge^pM\) 的基；\(p=0\) 时只有 \(e_{\varnothing}=1\)。对 \(p>n\)，有 \(\bigwedge^pM=0\)。特别地，在域上
\(\dim_k\bigwedge^p k^n=\binom np\)。
\end{theorem}
\begin{proof}
将各因子按 \(e_i\) 展开，再用多线性，任意纯楔积成为基向量楔积的有限线性组合。含重复指标的项为零，其余项由\eqref{b2:eq:alternating-permutation}化为带符号的递增指标项，故所列元素生成。当 \(p>n\) 时所有项都重复，得到零模。

为证明线性无关，希望对每个 \(e_I\) 构造一个只读取其系数的线性泛函，即在 \(e_I\) 上取一，在其他递增指标楔积上取零。先固定 \(1\le p\le n\) 及一个递增指标组 \(I\)。若 \(v_j=\sum_{a=1}^n x_{aj}e_a\)，每次从第 \(j\) 个输入选一个 \(I\) 中的基向量，按所选指标的排列符号求和，得到
\[
\Delta_I(v_1,\ldots,v_p)
=\sum_{\sigma\in S_p}\operatorname{sgn}(\sigma)
\prod_{j=1}^p x_{i_{\sigma(j)},j}.
\]
每个变量只出现一次，所以它多线性。若第 \(r,s\) 个输入相等，则把置换 \(\sigma\) 与 \(\sigma\circ(r\ s)\) 配对：二者不同，对应乘积相同而符号相反。每对的和为零，所以 \(\Delta_I\) 交替。这一成对消去在特征二中仍然有效。

泛性质给出线性映射 \(\widetilde\Delta_I:\bigwedge^pM\to R\)。对递增指标 \(J\)，当 \(I=J\) 时，求和中只有恒等置换对应的项为一，其余项均为零；当 \(I\ne J\) 时，每项都含有一个为零的坐标。因此 \(\widetilde\Delta_I(e_J)=1\) 当 \(I=J\)，否则为零。于是若 \(\sum_J a_Je_J=0\)，应用 \(\widetilde\Delta_I\) 便得 \(a_I=0\)。这证明线性无关，且没有除以任何整数。\(p=0\) 的结论由 \(\bigwedge^0M=R\) 直接得到。
\end{proof}

例如，在 \(R^3\) 中，二次外幂有基 \(e_1\wedge e_2,e_1\wedge e_3,e_2\wedge e_3\)。直接展开给出
\[
(e_1+2e_2)\wedge(e_2+e_3)
=e_1\wedge e_2+e_1\wedge e_3+2e_2\wedge e_3.
\]
系数 \(2\) 是否为零由所处的环决定，基的构造本身没有因此改变。

\subsection{分次交换律}
\label{b2:subsec:0902:02}

\begin{proposition}[分次交换律]\label{b2:prop:exterior-graded-commutativity}
设 \(R\) 为交换环，\(M\) 为 \(R\) 模，\(p,q\ge0\)。若 \(\xi\in\bigwedge^pM\)、\(\eta\in\bigwedge^qM\)，则
\begin{equation}\label{b2:eq:graded-wedge-sign}
\xi\wedge\eta=(-1)^{pq}\eta\wedge\xi.
\end{equation}
此外，每个奇数次齐次元素 \(\xi\) 都满足 \(\xi\wedge\xi=0\)，无须假设 \(2\) 可逆。
\end{proposition}
\begin{proof}
先取两个纯楔积。把前面的 \(p\) 个因子依次移过后面的 \(q\) 个因子，共交换 \(pq\) 次；每次由\cref{b2:prop:alternating-implies-skew}产生一个负号。这证明纯楔积情形，双线性再给出一般情形。

若 \(p\) 为奇数，分次交换律只直接给出 \(2\xi^2=0\)，在可能有二挠的环上不能据此推出平方为零。需要用外代数中重复因子为零的关系。写 \(\xi=\sum_{i=1}^t a_iw_i\)，其中 \(w_i\) 为 \(p\) 次纯楔积。每个 \(w_i\wedge w_i\) 都含重复因子，所以为零；不同指标的两项合为
\[
a_ia_j(w_i\wedge w_j+w_j\wedge w_i)=0
\]
，因为 \((-1)^{p^2}=-1\)。因此整个平方为零。
\end{proof}

若 \(2=0\) 于 \(R\)，分次交换律中的符号全部成为 \(1\)，外代数作为普通代数也是交换的，但仍保留每个一次元素的平方为零。例如
\[
\bigwedge(\mathbb F_2^2)
\cong\mathbb F_2[x,y]/(x^2,y^2),
\qquad e_1\longleftrightarrow x,\quad e_2\longleftrightarrow y.
\]
双方均由 \(1,x,y,xy\) 或相应楔积组成基，乘法关系一致，因而公式确为同构。它不同于无平方关系的对称代数 \(\mathbb F_2[x,y]\)。

外代数把每个一次元素的平方规定为零，而不只是把某一组基向量的平方规定为零。在线性映射的像中，\(u(x+y)^2=0\) 还要求两个交叉乘积之和为零。因此，即使 \(M\) 自由，仅检查基向量像的平方为零也不足以定义外代数上的同态。

\begin{proposition}[外代数的代数泛性质]\label{b2:prop:exterior-algebra-universal}
设 \(R\) 为交换环，\(M\) 为 \(R\) 模，\(A\) 为含单位元的结合 \(R\) 代数。若 \(R\) 线性映射 \(u:M\to A\) 满足 \(u(m)^2=0\) 对所有 \(m\in M\) 成立，则存在唯一 \(R\) 代数同态 \(\bigwedge M\to A\) 延伸 \(u\)。等价地，外代数是张量代数的商：
\[
\bigwedge M\cong T_R(M)/\langle m\otimes m\mid m\in M\rangle,
\]
分母表示这些二次元素生成的双边理想。
\end{proposition}
\begin{proof}
由假设，\(u(x)^2=u(y)^2=u(x+y)^2=0\)。由于 \(u\) 线性，展开第三个等式得
\[
0=(u(x)+u(y))^2
=u(x)^2+u(x)u(y)+u(y)u(x)+u(y)^2,
\]
因此 \(u(x)u(y)=-u(y)u(x)\)。乘积映射
\((m_1,\ldots,m_p)\mapsto u(m_1)\cdots u(m_p)\) 多线性；若两个输入相同，把其间的因子逐个交换，便出现一个 \(u(m)^2=0\)，所以它交替。逐次应用\cref{b2:prop:exterior-alternating-universal}得到各次外幂上的线性映射。首尾连接的公式保证它们合起来保持乘法和单位，故为代数同态。\(M\) 生成整个外代数，所以延伸唯一。

再由\cref{b2:thm:tensor-algebra-universal}，线性映射 \(M\to\bigwedge M\) 诱导满代数同态 \(T_R(M)\to\bigwedge M\)，并杀掉全部 \(m\otimes m\)，故经过上述商代数。反过来，商代数中的 \(M\) 的像都平方为零，由刚证明的泛性质得到反向代数同态。两复合在 \(M\) 与单位元上为恒等，因生成性而处处为恒等。
\end{proof}

\subsection{直和的外代数}
\label{b2:subsec:0902:03}

\begin{proposition}[直和的外幂]\label{b2:prop:exterior-direct-sum}
设 \(R\) 为交换环，\(M,N\) 为 \(R\) 模，\(p\ge0\)，则存在自然同构
\[
\bigoplus_{a+b=p}\left(\bigwedge^aM\otimes_R\bigwedge^bN\right)
\longrightarrow\bigwedge^p(M\oplus N),
\]
其在纯楔积上的公式为
\[
(m_1\wedge\cdots\wedge m_a)\otimes(n_1\wedge\cdots\wedge n_b)
\longmapsto
(m_1,0)\wedge\cdots\wedge(m_a,0)\wedge(0,n_1)\wedge\cdots\wedge(0,n_b).
\]
\end{proposition}
\begin{proof}
把 \(E=\bigwedge M\otimes_R\bigwedge N\) 按两侧次数之和分次。乘积原来的因子次序是 \(\xi,\eta,\xi',\eta'\)；要把 \(M\) 侧与 \(N\) 侧分别合在一起，就要让 \(N\) 侧次数为 \(b\) 的 \(\eta\) 越过 \(M\) 侧次数为 \(c\) 的 \(\xi'\)。在纯楔积上共发生 \(bc\) 次一次因子交换，所以应有符号 \((-1)^{bc}\)。据此采用带符号的分次张量积乘法，规定
\[
(\xi\otimes\eta)(\xi'\otimes\eta')
=(-1)^{bc}(\xi\wedge\xi')\otimes(\eta\wedge\eta'),
\qquad \deg\eta=b,\quad\deg\xi'=c.
\]
固定次数后此式分别线性且对两次张量关系平衡，故良定义。三因子结合律中，两种括号的符号指数分别为
\(bc+(b+d)e\) 与 \(de+b(c+e)\)，其中 \(d=\deg\eta'\)、\(e=\deg\xi''\)，二者相同；各外代数的结合律保证剩余因子也相同。单位是 \(1\otimes1\)。

映射 \(M\oplus N\to E\)、\((m,n)\mapsto m\otimes1+1\otimes n\) 的每个像都平方为零：两个同侧平方为零，两项交叉乘积分别为 \(m\otimes n\) 与 \(-m\otimes n\)。因此\cref{b2:prop:exterior-algebra-universal}给出分次代数同态 \(\alpha:\bigwedge(M\oplus N)\to E\)。

反方向，把 \(\xi\otimes\eta\) 送到两侧包含所诱导的楔积，得到线性映射 \(\beta:E\to\bigwedge(M\oplus N)\)。将中间的 \(\eta\) 与 \(\xi'\) 交换时，由\eqref{b2:eq:graded-wedge-sign}恰好得到 \((-1)^{bc}\)，所以 \(\beta\) 也是代数同态。两复合在所有 \(m\in M,n\in N\) 及单位元上都为恒等；这些元素生成双方代数，故 \(\alpha,\beta\) 互逆。取全次数 \(p\) 的分量，就得到命题中的直和同构。对 \(M,N\) 的线性映射，公式逐因子作用，故同构自然。
\end{proof}

一个 \(p\) 重楔积的输入来自 \(M\oplus N\)。将每个输入写成 \((m,n)=(m,0)+(0,n)\) 后按多线性展开，每一项恰有 \(a\) 个因子来自 \(M\)、\(b\) 个来自 \(N\)，其中 \(a+b=p\)。把 \(M\) 因子移到前面、\(N\) 因子移到后面，就落在上述直和的 \((a,b)\) 分量。因此直和公式是按两侧各取了多少个因子分类。特别地，
\[
\bigwedge^2(M\oplus N)
\cong\bigwedge^2M\oplus(M\otimes_RN)\oplus\bigwedge^2N.
\]
例如 \((m,0)\wedge(0,n)\) 属于中间项，而反向次序对应 \(-m\otimes n\)。这项负号已经由外代数的乘法决定，不是对直和分解额外指定的约定。若 \(M,N\) 分别具有 \(m,n\) 个基向量的有限自由模，比较基的数目便得到
\(\binom{m+n}{p}=\sum_{a+b=p}\binom ma\binom nb\)，其代数来源正是把一个递增指标组按两组基拆开。

\ProseParagraphBreak
在域上，纯楔积的非零性还能检测其因子是否独立。这里不需要先假定整个向量空间有限维：有限个因子只涉及一个有限维生成空间，再用直和公式比较它与母空间中的外幂。

\begin{proposition}\label{b2:prop:wedge-linear-independence}
设 \(k\) 为任意域，\(V\) 为 \(k\) 向量空间，\(p\ge1\)，\(v_1,\ldots,v_p\in V\)。则
\[
v_1\wedge\cdots\wedge v_p\ne0
\quad\Longleftrightarrow\quad
v_1,\ldots,v_p\text{ 线性无关}.
\]
\end{proposition}
\begin{proof}
若这些向量线性相关，取关系 \(\sum_{i=1}^p c_iv_i=0\) 中的非零系数 \(c_j\)，由域中可除性把 \(v_j\) 写成其余向量的线性组合。代入楔积并按多线性展开，每一项都含重复因子，因而为零。\(p=1\) 时这就是说 \(v_1=0\)，结论同样成立。

若这些向量线性无关，令 \(W=\operatorname{span}_k\{v_1,\ldots,v_p\}\)。它们构成有限维空间 \(W\) 的一组基，由\cref{b2:thm:exterior-power-basis}，\(v_1\wedge\cdots\wedge v_p\) 在 \(\bigwedge^pW\) 中是非零的基向量。由\cref{b2:thm:vector-basis-extension}将它们扩充为 \(V\) 的基，令 \(U\) 为其余基向量的线性生成空间，则 \(V=W\oplus U\)，即使 \(V\) 无限维也成立。\cref{b2:prop:exterior-direct-sum}给出
\[
\bigwedge^pV\cong
\bigoplus_{a+b=p}\bigl(\bigwedge^aW\otimes_k\bigwedge^bU\bigr).
\]
由该同构的纯楔积公式，\(W\hookrightarrow V\) 诱导的 \(\bigwedge^pW\rightarrow\bigwedge^pV\)，恰把 \(\bigwedge^pW\cong\bigwedge^pW\otimes_k k\) 送入 \((a,b)=(p,0)\) 的直和分量。因此它是单射，上述纯楔积在 \(\bigwedge^pV\) 中仍非零。
\end{proof}

在域上，独立向量还可以扩充成基；一般环上，非零楔积并不保证其因子能够扩充成母模的一组基。例如在 \(\mathbb Z^2\) 中，
\[
(2e_1)\wedge e_2=2(e_1\wedge e_2)\ne0,
\]
因为 \(e_1\wedge e_2\) 是自由 \(\mathbb Z\) 模 \(\bigwedge^2\mathbb Z^2\) 的基向量。但 \(2e_1,e_2\) 只能生成第一坐标为偶数的子模，不能生成整个 \(\mathbb Z^2\)。它们虽独立，楔积中的系数 \(2\) 却不是单位；“系数非零”与“系数可逆”的差别将在行列式的可逆性判据中再次出现。

% !TeX root = ../../Book2.tex
% 纲要标题：### 9.3 行列式的内在定义
% 纲要位置：计划纲要.md，第 1080 行。

\section{行列式的内在定义}
\label{b2:sec:0903}

一个 \(n\) 维空间中的 \(n\) 重楔积只剩下一条线，所以线性变换在这条线上只能表现为乘以一个标量。这个标量就是行列式。以最高外幂为起点，乘法性来自函子复合，子式公式来自较低次外幂的坐标。本节同时允许任意交换环上的有限自由模，不以域上的除法或特征值作为前提。

% 纲要要点（供目录核对）：
%
% * 最高外幂与行列式线
% * 行列式的乘法性
% * Cauchy–Binet 公式的外代数解释
%

\subsection{最高外幂与行列式线}
\label{b2:subsec:0903:01}

\begin{definition}\label{b2:def:determinant-line}
设 \(R\) 为交换环，\(M\) 为有 \(n\) 个基向量的有限自由 \(R\) 模。它的最高外幂
\[
\det M\coloneqq\bigwedge^nM
\]
是秩一自由 \(R\) 模，称为 \(M\) 的\term[hanglieshi-xian]{行列式线}[determinant line]。\(M=0\) 且取空基时，约定 \(\det M=R\)。
\end{definition}
\symidx{\(\det M=\bigwedge^nM\)：有限自由模的行列式线}

由\cref{b2:thm:exterior-power-basis}，任一有序基 \(e_1,\ldots,e_n\) 给出行列式线的基向量 \(\omega=e_1\wedge\cdots\wedge e_n\)。对于线性自同态 \(f:M\to M\)，存在唯一 \(a\in R\) 使
\begin{equation}\label{b2:eq:intrinsic-determinant}
\left(\bigwedge^n f\right)(\omega)=a\omega.
\end{equation}
这个 \(a\) 不依赖 \(\omega\) 的选择。事实上，同一自由模的另一个单元素基必为 \(\omega'=u\omega\)，其中 \(u\) 为单位：若 \(\omega=v\omega'\)，则由基坐标唯一性得 \(uv=1\)。现在
\(\bigwedge^n f(\omega')=u a\omega=a\omega'\)，使用的正是 \(R\) 的交换性。

\ProseParagraphBreak
据此定义 \(f\) 的\term[hanglieshi]{行列式}[determinant]为 \(\det f=a\)。对 \(n\times n\) 矩阵 \(A=(a_{ij})\)，将其看成 \(R^n\) 上的线性映射并记其行列式为 \(\det A\)。把列向量 \(f(e_j)=\sum_i a_{ij}e_i\) 代入\eqref{b2:eq:intrinsic-determinant}，多线性展开后只有没有重复指标的项保留，故
\begin{equation}\label{b2:eq:determinant-leibniz}
\det A=\sum_{\sigma\in S_n}\operatorname{sgn}(\sigma)
\prod_{j=1}^n a_{\sigma(j),j}.
\end{equation}
这与熟悉的 Leibniz 公式相同；本节的定义没有改变原来的行列式，而是解释了它为何具有坐标无关的含义。空矩阵的行列式为 \(1\)，与空楔积的约定一致。
\symidx{\(\det f,\det A\)：线性自同态及矩阵的行列式}

\ProseParagraphBreak
公式立即表明行列式对各列分别线性，若两列相同则为零，交换两列变号。对转置矩阵，用置换 \(\sigma^{-1}\) 重排求和，并利用标量乘法交换，得到
\begin{equation}\label{b2:eq:determinant-transpose}
\det(A^{\mathsf t})=\det A.
\end{equation}
因此相同性质也适用于行。上三角矩阵的行列式等于对角元之积：非恒等置换必有某个 \(\sigma(j)>j\)，相应乘积含下三角的零元，只有恒等置换可能贡献。

\subsection{行列式的乘法性}
\label{b2:subsec:0903:02}

\begin{theorem}[行列式的乘法性]\label{b2:thm:determinant-multiplicativity}
设 \(R\) 为交换环，\(M\) 为有限自由 \(R\) 模，\(f,g:M\to M\) 为 \(R\) 线性自同态，则
\[
\det(gf)=\det(g)\det(f),\qquad
\det(\operatorname{id})=1.
\]
特别地，\(n\times n\) 矩阵满足 \(\det(BA)=\det B\det A\)。若 \(P\) 可逆，则
\[
\det(P^{-1}AP)=\det A.
\]
\end{theorem}
\begin{proof}
对行列式线的基 \(\omega\)，由\eqref{b2:eq:exterior-functoriality}，
\[
\begin{aligned}
\left(\bigwedge^n(gf)\right)(\omega)
&=\left(\bigwedge^n g\right)\left(\bigwedge^n f\right)(\omega)\\
&=\det(f)\det(g)\omega.
\end{aligned}
\]
基系数唯一，得到乘法公式；恒等映射保持 \(\omega\)，得到第二式。若 \(P\) 可逆，乘法性给出 \(\det(P^{-1})\det P=1\)，所以相似矩阵公式由三次乘法直接推出。
\end{proof}

在域上，\(\det f\ne0\) 等价于 \(f\) 可逆：若像中的 \(n\) 个向量线性相关，其楔积为零；若线性无关，它们构成一组基，从而 \(f\) 可逆。这里“相关则楔积为零”由把一个向量写成其他向量的线性组合得到；反方向把无关组扩充为基，再用\cref{b2:thm:exterior-power-basis}可知其楔积非零。

\ProseParagraphBreak
在一般交换环上，非零行列式不保证可逆。例如 \(\mathbb Z\) 上乘二的行列式为 \(2\ne0\)，但乘二不满射。正确条件是行列式为单位；必要性已经由乘法性得到，充分性将在\cref{b2:prop:adjugate-identity}中用伴随矩阵证明。单射、满射与可逆也不宜再用域上的秩直觉一概替代。

\subsection{Cauchy–Binet 公式的外代数解释}
\label{b2:subsec:0903:03}

设 \(A:R^n\to R^m\) 的矩阵仍记为 \(A\)。对递增指标集 \(I\subseteq\{1,\ldots,m\}\)、\(J\subseteq\{1,\ldots,n\}\)，且 \(|I|=|J|=p\)，记 \(A_{I,J}\) 为对应的 \(p\times p\) 子矩阵；它的行列式称为一个 \(p\) 阶\term[zishi]{子式}[minor]。基向量楔积的展开与\eqref{b2:eq:determinant-leibniz}给出
\begin{equation}\label{b2:eq:exterior-minor-action}
\left(\bigwedge^p A\right)(e_J)
=\sum_{|I|=p}\det(A_{I,J})e_I.
\end{equation}
这里目标与来源分别使用各自的标准基；每个系数都是选定 \(p\) 列后，对相应 \(p\) 行坐标所作的交替求和。因此，较低次外幂把所有同阶子式同时组织成一个线性映射。

\begin{theorem}[Cauchy–Binet]\label{b2:thm:cauchy-binet}
设 \(R\) 为交换环，\(m,n,r\ge0\)，\(A\in M_{m\times n}(R)\)、\(B\in M_{n\times r}(R)\)，\(p\) 为满足 \(0\le p\le\min(m,r)\) 的整数。对大小均为 \(p\) 的递增指标集 \(I\subseteq\{1,\ldots,m\}\)、\(J\subseteq\{1,\ldots,r\}\)，有
\begin{equation}\label{b2:eq:cauchy-binet}
\det\bigl((AB)_{I,J}\bigr)
=\sum_{\substack{K\subseteq\{1,\ldots,n\}\\|K|=p}}
\det(A_{I,K})\det(B_{K,J}).
\end{equation}
若 \(p>n\)，右端为空和，等于零。特别地，\(A\) 为 \(p\times n\)、\(B\) 为 \(n\times p\) 矩阵时，此式给出 \(\det(AB)\) 的全部最大子式展开。
\end{theorem}
\begin{proof}
先对 \(\bigwedge^pB\) 使用\eqref{b2:eq:exterior-minor-action}，再对每个结果使用 \(\bigwedge^pA\)，得到
\[
\left(\bigwedge^pA\right)\left(\bigwedge^pB\right)(e_J)
=\sum_I\left(\sum_K\det(A_{I,K})\det(B_{K,J})\right)e_I.
\]
另一方面，函子性说明左边等于 \(\bigwedge^p(AB)(e_J)\)，其 \(e_I\) 系数是 \(\det\bigl((AB)_{I,J}\bigr)\)。外幂基的坐标唯一性给出公式。若 \(p>n\)，中间模 \(\bigwedge^p R^n\) 为零，复合映射为零，故全部这些子式都为零。
\end{proof}

这一公式称为\term[Cauchy-Binet-gongshi]{Cauchy–Binet 公式}[Cauchy--Binet formula]。例如取
\[
A=\begin{pmatrix}1&0&1\\0&1&1\end{pmatrix},\qquad
B=\begin{pmatrix}1&2\\0&1\\1&0\end{pmatrix}.
\]
直接相乘得到 \(AB=\begin{pmatrix}2&2\\1&1\end{pmatrix}\)，行列式为零。按列指标 \(12,13,23\)，\(A\) 的三个二阶子式为 \(1,1,-1\)，\(B\) 对应行指标的子式为 \(1,-2,-1\)，故 Cauchy–Binet 求和为 \(1-2+1=0\)。这个算例同时核对了指标次序和符号。

\ProseParagraphBreak
当 \(M,N\) 分别具有 \(m,n\) 个基向量时，\cref{b2:prop:exterior-direct-sum}在次数 \(m+n\) 只剩一项，给出
\[
\det(M\oplus N)\cong\det M\otimes_R\det N.
\]
在这个同构下，\(f\oplus g\) 的最高外幂是两条行列式线上的作用之张量积，所以块对角矩阵的行列式为两个块的行列式之积。这也是直和外代数公式在最高次数的直接应用。

% !TeX root = ../../Book2.tex
% 纲要标题：### 9.4 外代数的应用
% 纲要位置：计划纲要.md，第 1086 行。

\section{外代数的应用}
\label{b2:sec:0904}

外代数既解释行列式，也提供记录子空间与构造复形的方法。本节先从子式推出伴随矩阵公式，并在任意交换环上证明 Cayley–Hamilton 定理；然后在域上用一个四维例子说明 Plücker 关系；最后给出 Koszul 微分的构造。三个应用分别用到最高外幂、二次外幂和整个分次外代数。

% 纲要要点（供目录核对）：
%
% * 余子式、伴随矩阵与 Cayley–Hamilton；系数递推和相邻项抵消
% * Plücker 坐标的初步例子；非零纯楔积恢复子空间、四维二次关系与特征二
% * Koszul 复形的外代数背景；收缩由分次乘法规则产生，反交换关系完整证明
% * 三变量自然关系与关系之间的相容性；复合为零不等于正合，自由消解的含义
% * 含单位序列的 Koszul 复形正合，构造及证明放在正文；正则序列理论归第三卷
%
% ---
%

\subsection{余子式、伴随矩阵与 Cayley–Hamilton 定理}
\label{b2:subsec:0904:01}

设 \(R\) 为交换环，\(A=(a_{ij})\in M_n(R)\)、\(n\ge1\)。删去第 \(i\) 行、第 \(j\) 列所得子矩阵记为 \(A^{\widehat i,\widehat j}\)。数
\[
C_{ij}=(-1)^{i+j}\det A^{\widehat i,\widehat j}
\]
称为对应的\term[daishu-yuzishi]{代数余子式}[cofactor]。当 \(n=1\) 时，删去后是空矩阵，行列式取 \(1\)。定义
\[
\operatorname{adj}(A)=(C_{ji})_{i,j}
\]
为 \(A\) 的伴随矩阵（adjugate matrix）\termidx[ban-sui-juzhen]{伴随矩阵}。这是代数余子式矩阵的转置，与第四章由首一多项式构造的友矩阵（companion matrix）是不同概念。
\symidx{\(\operatorname{adj}(A)\)：矩阵 \(A\) 的伴随矩阵}

\begin{proposition}[余子式展开与伴随矩阵]\label{b2:prop:adjugate-identity}
设 \(R\) 为交换环，\(A=(a_{ij})\in M_n(R)\)、\(n\ge1\)。令 \(C_{ij}=(-1)^{i+j}\det A^{\widehat i,\widehat j}\) 为代数余子式，\(\operatorname{adj}(A)=(C_{ji})_{i,j}\) 为伴随矩阵，其中 \(n=1\) 时空矩阵的行列式取 \(1\)。固定任意行 \(i\) 或列 \(j\)，有
\[
\det A=\sum_{k=1}^n a_{ik}C_{ik}
=\sum_{k=1}^n a_{kj}C_{kj}.
\]
进而
\begin{equation}\label{b2:eq:adjugate-identity}
A\operatorname{adj}(A)=\operatorname{adj}(A)A=(\det A)I_n.
\end{equation}
所以 \(A\) 可逆当且仅当 \(\det A\) 是 \(R\) 的单位，此时
\(A^{-1}=(\det A)^{-1}\operatorname{adj}(A)\)。
\end{proposition}
\begin{proof}
在\eqref{b2:eq:determinant-leibniz}中，按某一固定行所使用的列位置分组。含 \(a_{ik}\) 的项，其余因子的排列恰好遍历删去该行、该列后的置换。把第 \(i\) 行和第 \(k\) 列分别移到首位，共产生 \(i-1+k-1\) 个换位，故这一组的符号因子为 \((-1)^{i+k}\)。组内求和于是等于 \(a_{ik}C_{ik}\)，得到行展开。对转置矩阵使用行展开，再用\eqref{b2:eq:determinant-transpose}，得到列展开。

乘积 \(A\operatorname{adj}(A)\) 的 \((i,j)\) 项为 \(\sum_k a_{ik}C_{jk}\)。当 \(i=j\) 时是 \(\det A\)。当 \(i\ne j\) 时，我们仍要把这个和解释成一次行列式展开：把 \(A\) 的第 \(j\) 行替换为第 \(i\) 行，沿第 \(j\) 行展开时，代数余子式仍为 \(C_{jk}\)，行中的系数则变为 \(a_{ik}\)。所得行列式正是上述和，但矩阵有两行相同，故为零。这证明第一个乘积公式。第二个乘积公式通过替换第 \(i\) 列为第 \(j\) 列并作列展开得到。

若行列式为单位，式\eqref{b2:eq:adjugate-identity}给出双侧逆。若 \(A\) 可逆，则\cref{b2:thm:determinant-multiplicativity}给出 \(\det A\det A^{-1}=1\)，所以行列式为单位。
\end{proof}

下面通过比较多项式的系数重新证明 Cayley–Hamilton 定理\termidx[Cayley-Hamilton-dingli]{Cayley--Hamilton 定理}。\cref{b2:prop:operator-characteristic-minimal}利用域上的多项式环模分类，证明特征多项式是最小多项式的倍数，因而零化算子；这里直接使用行列式与伴随矩阵恒等式，把系数推广到任意交换环。

\begin{theorem}[Cayley–Hamilton]\label{b2:thm:exterior-cayley-hamilton}
若 \(R\) 为交换环、\(A\in M_n(R)\)、\(n\ge1\)，令
\[
\chi_A(t)=\det(tI_n-A)=t^n+c_{n-1}t^{n-1}+\cdots+c_0.
\]
则 \(\chi_A(A)=A^n+c_{n-1}A^{n-1}+\cdots+c_0I_n=0\)。
\end{theorem}
\begin{proof}
在交换环 \(R[t]\) 上应用\cref{b2:prop:adjugate-identity}，得到
\[
(tI_n-A)\operatorname{adj}(tI_n-A)=\chi_A(t)I_n.
\]
对标量多项式 \(\chi_A(t)\)，\(\chi_A(A)\) 当然有意义。但若要在左端按普通多项式代入的规则令 \(t=A\)，还须先核对矩阵系数与 \(A\) 交换：在矩阵系数多项式中，\(t\) 与所有系数交换，换成 \(A\) 后却不能未经证明地保留这一性质。因此这里保留每个矩阵因子的次序，改由系数关系消去伴随矩阵中的未知系数。

伴随矩阵的每个条目是 \(n-1\) 阶子式，所以次数至多 \(n-1\)。将它按 \(t\) 展开为
\[
\operatorname{adj}(tI_n-A)=B_0+B_1t+\cdots+B_{n-1}t^{n-1},
\qquad B_j\in M_n(R).
\]
比较 \(t\) 的各次系数，依次得到
\[
-AB_0=c_0I_n,\qquad
B_{j-1}-AB_j=c_jI_n\ (1\le j<n),\qquad B_{n-1}=I_n.
\]
这些关系把 \(c_jI_n\) 写成相邻两个 \(B_j\) 的组合，不必逐个求出 \(B_j\)。将中间第 \(j\) 个等式在左侧乘以 \(A^j\)，最后一式在左侧乘以 \(A^n\)，再连同首式求和，相邻的 \(B_j\) 项便有相同的左侧幂次而符号相反。右端为 \(\chi_A(A)\)，左端为
\[
-AB_0+\sum_{j=1}^{n-1}(A^jB_{j-1}-A^{j+1}B_j)
+A^nB_{n-1}=0.
\]
相邻项逐一抵消，故 \(\chi_A(A)=0\)。这一计算保持每个矩阵系数的位置，没有使用矩阵环中任意元素交换的假设。
\end{proof}

例如，对 \(A=\begin{pmatrix}a&b\\c&d\end{pmatrix}\)，有
\[
\chi_A(t)=t^2-(a+d)t+(ad-bc),
\qquad A^2-(a+d)A+(ad-bc)I_2=0.
\]
这些等式在 \(\mathbb Z\)、多项式环以及含幂零元的交换环上都成立。证明中用到的交换性属于标量环，而不是所有 \(2\times2\) 矩阵之间的交换性。

\subsection{Plücker 坐标}
\label{b2:subsec:0904:02}

本小节改设 \(k\) 为域、\(V=k^n\)。对 \(r\) 维子空间 \(W\subseteq V\)，选一组有序基 \(w_1,\ldots,w_r\)，写
\[
w_1\wedge\cdots\wedge w_r=\sum_{|I|=r}p_Ie_I.
\]
这些 \(p_I\) 是以 \(w_j\) 为列的矩阵的 \(r\) 阶子式。若换一组基，整个楔积乘以换基矩阵的非零行列式。因此，把所有 \(p_I\) 同时乘以同一个非零标量视为同一组坐标，就得到只依赖 \(W\) 的\term[Plucker-zuobiao]{Plücker 坐标}[Plücker coordinates]。这种只确定到共同非零倍数的坐标称为齐次坐标；子空间 \(W\) 确定的是 \(\bigwedge^rV\) 中的一条直线 \(k(w_1\wedge\cdots\wedge w_r)\)，而不是这条直线上的某个指定非零向量。

\ProseParagraphBreak
这些坐标确实能够辨认子空间。对非零纯楔积 \(\omega=w_1\wedge\cdots\wedge w_r\)，有
\begin{equation}\label{b2:eq:plucker-recovers-plane}
W=\{\,v\in V\mid v\wedge\omega=0\,\}.
\end{equation}
若 \(v\in W\)，按 \(w_j\) 展开即得楔积为零。若 \(v\notin W\)，则 \(v,w_1,\ldots,w_r\) 线性无关，\cref{b2:prop:wedge-linear-independence}说明它们的楔积非零。因此两个非零纯楔积只相差标量时，恢复出的子空间相同。这里的非零与纯楔积条件都不可省略；对一般外张量，右端不能无条件解释为一个 \(r\) 维子空间。\cref{b2:ex:09-plucker-plane-recovery}给出了由具体坐标恢复平面的计算。

\ProseParagraphBreak
还需判定哪些坐标真的来自一个子空间。例如，若一个 \(4\times2\) 矩阵的两列为 \(u,v\)，它的六个二阶子式同时来自这两列，因而必须满足彼此的相容关系。下面的二次式正是这种约束。

\begin{proposition}[四维空间中二重楔积的判据]\label{b2:prop:plucker-four-dimensional}
设 \(k\) 为任意域，\(e_1,\ldots,e_4\) 为 \(k^4\) 的标准基，并设
\[
\omega=\sum_{1\le i<j\le4}p_{ij}e_i\wedge e_j\in\bigwedge^2k^4.
\]
它可以写为 \(u\wedge v\)，当且仅当
\begin{equation}\label{b2:eq:plucker-quadric}
p_{12}p_{34}-p_{13}p_{24}+p_{14}p_{23}=0.
\end{equation}
此结论对任意特征的域成立。\(\omega\ne0\) 时，\(u,v\) 自动线性无关。
\end{proposition}
\begin{proof}
若 \(u=\sum u_ie_i\)、\(v=\sum v_ie_i\)，则 \(p_{ij}=u_iv_j-u_jv_i\)。代入\eqref{b2:eq:plucker-quadric}，得到
\[
\begin{aligned}
&(u_1v_2-u_2v_1)(u_3v_4-u_4v_3)\\
&\quad-(u_1v_3-u_3v_1)(u_2v_4-u_4v_2)\\
&\quad+(u_1v_4-u_4v_1)(u_2v_3-u_3v_2)=0.
\end{aligned}
\]
展开后的十二项按相同的 \(u_iu_jv_av_b\) 配对，符号相反，所以为零。

反过来，\(\omega=0\) 时取 \(u=0\)。若 \(\omega\ne0\)，至少一个 \(p_{ij}\ne0\)。重排基可使它成为 \(p_{12}\)：对相邻交换 \((1\ 2),(2\ 3),(3\ 4)\)，直接将下标重排并用 \(p_{ji}=-p_{ij}\)，可见关系左边都变为其负元，所以关系在重排后仍成立。令 \(a=p_{12}\ne0\)。我们先把待求的两列 \(u,v\) 的前两行选为 \((1,0)\)、\((0,a)\)，使 \(12\) 子式已经等于 \(a\)：
\[
\begin{pmatrix}
1&0\\
0&a\\
u_3&v_3\\
u_4&v_4
\end{pmatrix}.
\]
再匹配四个包含第 \(1\) 或第 \(2\) 行的子式，就得到
\[
v_3=p_{13},\qquad v_4=p_{14},\qquad
-a u_3=p_{23},\qquad -a u_4=p_{24}.
\]
前两行使这四个子式成为关于未知系数的线性方程。由于 \(a\ne0\)，它们唯一确定 \(u_3,u_4,v_3,v_4\)，即
\[
u=e_1-\frac{p_{23}}a e_3-\frac{p_{24}}a e_4,
\qquad v=a e_2+p_{13}e_3+p_{14}e_4.
\]
展开 \(u\wedge v\)，\(12,13,14,23,24\) 五个坐标分别就是给定的 \(p_{ij}\)，而 \(34\) 坐标为
\[
\frac{p_{13}p_{24}-p_{14}p_{23}}a=p_{34},
\]
最后一步恰用\eqref{b2:eq:plucker-quadric}：五个坐标已由所选两列实现，二次关系保证剩下的第六个坐标也匹配。故 \(u\wedge v=\omega\)。若 \(u,v\) 相关，其楔积为零，所以非零情形下它们独立。
\end{proof}

外幂中能写成一个纯楔积的元素称为\term[kefenjie-zhangliang]{可分解外张量}[decomposable exterior tensor]，也简称可分解张量。这里的“可分解”专指能写成一个楔积，区别于张量积中能写成一个纯张量。上面的关系表明，在 \(\bigwedge^2k^4\) 中，可分解外张量的坐标并不是任意六个标量。例如 \(e_1\wedge e_2+e_3\wedge e_4\) 的关系左端为 \(1\)，所以在任意域上都不可分解。一般 \(r\) 维子空间的 Plücker 坐标也满足一族二次关系；这里的 \(r=2,n=4\) 是只需一个非平凡二次方程的最小例子。

\ProseParagraphBreak
另一方面，简记 \(e_{ij}=e_i\wedge e_j\)、\(e_{1234}=e_1\wedge e_2\wedge e_3\wedge e_4\)。展开 \(\omega\wedge\omega\) 时，只有指标互不重复的交叉项可能非零。两个次数二的因子交换时产生的符号为 \((-1)^{2\cdot2}=1\)，所以 \(e_{12}\wedge e_{34}\) 与 \(e_{34}\wedge e_{12}\) 同号，都等于 \(e_{1234}\)。每一对这样的交叉项因而贡献两倍，得到
\[
\omega\wedge\omega
=2(p_{12}p_{34}-p_{13}p_{24}+p_{14}p_{23})
e_1\wedge e_2\wedge e_3\wedge e_4.
\]
特征不为二时，\(\omega\wedge\omega=0\) 与 Plücker 关系等价；特征二时，每个这样的平方都是零，包括刚才的不可分解张量。因而在特征二中仍应使用坐标关系本身来判定，不能把平方为零当作充分条件。

\subsection{Koszul 复形的外代数背景}
\label{b2:subsec:0904:03}

重新回到任意交换环 \(R\)。给定整数 \(n\ge1\) 与 \(a_1,\ldots,a_n\in R\)，以 \(e_1,\ldots,e_n\) 为 \(R^n\) 的标准基。生成元之间的关系由映射
\[
\lambda:R^n\longrightarrow R,\qquad e_i\longmapsto a_i
\]
的核记录：\(\sum_i r_ie_i\) 在核中，恰好表示 \(\sum_i r_i a_i=0\)。交换性立即给出关系 \(a_ie_j-a_je_i\)，因为它的像是 \(a_i a_j-a_j a_i=0\)。二次外幂中的 \(e_i\wedge e_j\) 正好记录这一对指标，更高外幂则用来组织这些关系之间的相容性。

\ProseParagraphBreak
对任意 \(R\) 模 \(M\)，为了把线性映射 \(\lambda:M\to R\) 延伸到整个外代数，我们要求它在一次分量上等于 \(\lambda\)，在标量上为零，并遵循分次 Leibniz 规则：作用于一个楔积时，可以作用于前一因子，也可以越过前一因子后作用于后一因子；越过次数 \(p\) 的因子时带上符号 \((-1)^p\)。次数为 \(-1\) 表示把 \(\bigwedge^pM\) 映到 \(\bigwedge^{p-1}M\)。若这样的算子连续作用两次总为零，即平方为零，就称为微分。本小节将所有负次外幂约定为零，使零次上的微分也有明确的目标。

\begin{proposition}[收缩微分]\label{b2:prop:exterior-contraction}
设 \(R\) 为交换含幺环，\(M\) 为幺性 \(R\) 模，\(\lambda:M\to R\) 为 \(R\) 线性映射，并约定 \(\bigwedge^pM=0\) 对 \(p<0\) 成立。存在唯一次数为 \(-1\) 的 \(R\) 线性分次导子
\[
\iota_\lambda:\bigwedge^pM\longrightarrow\bigwedge^{p-1}M,
\]
在 \(R=\bigwedge^0M\) 上取零，在 \(M=\bigwedge^1M\) 上等于 \(\lambda\)，且对 \(\xi\in\bigwedge^pM\)、\(\eta\in\bigwedge^qM\) 满足分次 Leibniz 规则
\[
\iota_\lambda(\xi\wedge\eta)
=\iota_\lambda(\xi)\wedge\eta+(-1)^p\xi\wedge\iota_\lambda(\eta).
\]
它在纯楔积上的公式为
\begin{equation}\label{b2:eq:contraction-formula}
\iota_\lambda(m_1\wedge\cdots\wedge m_p)
=\sum_{j=1}^p(-1)^{j-1}\lambda(m_j)
m_1\wedge\cdots\wedge\widehat{m_j}\wedge\cdots\wedge m_p.
\end{equation}
这里 \(p\ge1\)，帽号表示省去该因子。这个算子称为由 \(\lambda\) 确定的\term[shousuo]{收缩}[contraction]，并满足 \(\iota_\lambda^2=0\)。
\end{proposition}
\begin{proof}
若 \(D\) 是满足陈述中条件的分次导子，先把第一个因子分出，Leibniz 规则给出
\[
D(m_1\wedge\cdots\wedge m_p)
=\lambda(m_1)m_2\wedge\cdots\wedge m_p
-m_1\wedge D(m_2\wedge\cdots\wedge m_p).
\]
从 \(p=1\) 开始归纳，便强迫得到\eqref{b2:eq:contraction-formula}。纯楔积生成各次外幂，而零次上的作用已经指定，所以至多存在一个这样的算子。

反过来，式\eqref{b2:eq:contraction-formula}的右端对输入多线性。若 \(m_r=m_s\)、\(r<s\)，除删除第 \(r\) 或第 \(s\) 个因子的项以外，其余项仍有重复因子，均为零。删去第 \(r\) 个因子后，把原第 \(s\) 个因子移回第 \(r\) 个位置，共交换 \(s-r-1\) 次，所以该项的总符号是 \((-1)^{r-1+s-r-1}=(-1)^{s-2}\)，与删去第 \(s\) 个因子的项符号相反。二者系数中的 \(\lambda(m_r),\lambda(m_s)\) 相同，故抵消。这证明右端交替，外幂泛性质给出各次上的线性算子。把零次上的作用取为零，就得到次数为 \(-1\) 的 \(\iota_\lambda\)，其在一次分量上的作用是 \(\lambda\)。

当 \(p\) 或 \(q\) 为零时，所示乘法公式由 \(R\) 线性与标量上的零作用直接得到。对两个正次纯楔积的乘积，删除位置或者在前 \(p\) 项，或者在后 \(q\) 项。前者给出第一项，后者相对从 \(\eta\) 开始计数多出 \(p\)，产生 \((-1)^p\)，得到分次 Leibniz 规则。由双线性推广到所有齐次元素。

最后，两次收缩就是选两个不同位置，各施加一次 \(\lambda\) 并删去相应因子。每对位置有两种删除顺序，它们留下相同的楔积与相同的标量，只在符号上相反。具体说，固定 \(i<j\)，先删 \(i\) 再删原第 \(j\) 个因子的符号为 \((-1)^{i+j-3}\)，先删 \(j\) 再删 \(i\) 的符号为 \((-1)^{i+j-2}\)。两项的剩余因子次序相同，标量系数分别为 \(\lambda(m_i)\lambda(m_j)\) 及其反序，由交换性相等，故相互抵消。所有项成对消去，得到 \(\iota_\lambda^2=0\)。
\end{proof}
\symidx{\(\iota_\lambda\)：外代数上的收缩算子}

\begin{corollary}\label{b2:cor:contractions-anticommute}
设 \(R\) 为交换含幺环，\(M\) 为幺性 \(R\) 模，\(\lambda,\mu:M\rightarrow R\) 为 \(R\) 线性映射。由它们确定的外代数收缩算子满足
\[
\iota_\lambda\iota_\mu+\iota_\mu\iota_\lambda=0.
\]
这个反交换关系对任意特征成立。
\end{corollary}
\begin{proof}
由\eqref{b2:eq:contraction-formula}，\(\iota_{\lambda+\mu}\) 与 \(\iota_\lambda+\iota_\mu\) 在每个纯楔积上相同，在零次部分上也都为零，所以
\(\iota_{\lambda+\mu}=\iota_\lambda+\iota_\mu\)。对 \(\lambda+\mu\) 应用\cref{b2:prop:exterior-contraction}，得到
\[
0=\iota_{\lambda+\mu}^2
=\iota_\lambda^2+\iota_\lambda\iota_\mu
+\iota_\mu\iota_\lambda+\iota_\mu^2.
\]
两个平方项已为零，剩下的就是所求关系。这里只展开算子乘积并消去已知零项，没有除以二，因此特征二时也成立。
\end{proof}

令 \(n\ge1\)、\(M=R^n\)，选基 \(e_1,\ldots,e_n\)，并指定 \(\lambda(e_i)=a_i\in R\)。取 \(K_p=\bigwedge^pR^n\)、\(\mathrm{d}_p=\iota_\lambda\)，得到
\[
0\longrightarrow K_n\xrightarrow{\mathrm{d}_n}K_{n-1}\longrightarrow\cdots
\longrightarrow K_1\xrightarrow{\mathrm{d}_1}K_0\longrightarrow0.
\]
一列 \(R\) 模与 \(R\) 线性箭头若满足每两个相邻箭头的复合为零，就称为一个\term[lianfuxing]{链复形}[chain complex]；此处的条件 \(\mathrm{d}_{p-1}\mathrm{d}_p=0\) 已由\cref{b2:prop:exterior-contraction}证明。上面由 \(a_1,\ldots,a_n\in R\) 构造的链复形称为序列 \(a_1,\ldots,a_n\) 的\term[Koszul-fuxing]{Koszul 复形}[Koszul complex]，记为 \(K_\bullet(a_1,\ldots,a_n;R)\)。\symidx{\(K_\bullet(a_1,\ldots,a_n;R)\)：Koszul 复形}它在低次上的公式是
\[
\mathrm{d}_1(e_i)=a_i,\qquad \mathrm{d}_2(e_i\wedge e_j)=a_ie_j-a_je_i.
\]
\(K_1=R^n\) 记录这组理想生成元，\(K_2\) 记录由交换性自动产生的两两关系。因此 \(\operatorname{im}\mathrm{d}_1=(a_1,\ldots,a_n)\)，它的余核就是 \(R/(a_1,\ldots,a_n)\)；把这个商接到末端，是由 \(\mathrm{d}_1\) 的像决定的。得到增广复形
\[
0\longrightarrow K_n\longrightarrow\cdots\longrightarrow K_1
\xrightarrow{\mathrm{d}_1}R\longrightarrow R/(a_1,\ldots,a_n)\longrightarrow0.
\]
例如，当 \(n=2\) 时，未增广的复形为
\[
0\longrightarrow R\xrightarrow{\binom{-a_2}{a_1}}R^2
\xrightarrow{\begin{psmallmatrix}a_1&a_2\end{psmallmatrix}}R\longrightarrow0.
\]
相邻矩阵的乘积为 \(-a_1a_2+a_2a_1=0\)，\(\mathrm{d}_2\) 记录了最直接的两两关系。

\ProseParagraphBreak
当 \(n=3\)，写 \(a_1=a,a_2=b,a_3=c\)，并简记 \(e_{ij}=e_i\wedge e_j\)、\(e_{123}=e_1\wedge e_2\wedge e_3\)。三条两两关系为 \(ae_2-be_1\)、\(ae_3-ce_1\)、\(be_3-ce_2\)，而三次外幂给出
\[
\mathrm{d}_3(e_{123})=a e_{23}-b e_{13}+c e_{12}.
\]
它说明这三条关系本身也有相容关系，因为
\[
\begin{aligned}
\mathrm{d}_2\mathrm{d}_3(e_{123})
&=a(be_3-ce_2)-b(ae_3-ce_1)+c(ae_2-be_1)\\
&=0.
\end{aligned}
\]
于是三次外幂记录的是二次外幂所给关系之间的关系，更高次外幂继续以相同方式组织这些自然关系的相容性。

\ProseParagraphBreak
这些自然关系未必生成全部关系。取任意域 \(k\)，令 \(R=k[x]\)、\(a_1=a_2=x\)。此时 \(\mathrm{d}_1(u,v)=x(u+v)\)，\(\mathrm{d}_2(t)=(-xt,xt)\)。因为 \(R\) 是整环且 \(x\ne0\)，有
\[
\ker\mathrm{d}_1=R(-1,1),\qquad
\operatorname{im}\mathrm{d}_2=xR(-1,1)\subsetneq R(-1,1).
\]
严格包含来自 \(1\notin xR\)：关系 \((-1,1)\) 在核中，却不是 \(\mathrm{d}_2\) 的像。因此序列 \(x,x\) 的 Koszul 复形在 \(R^2\) 处不正合。

\ProseParagraphBreak
复合为零给出 \(\operatorname{im}\mathrm{d}_p\subseteq\ker\mathrm{d}_{p-1}\)；在相应位置正合则要求二者相等。增广复形在 \(R\) 及商模处正合，但在其他位置未必正合。一条增广链复形若在每个位置正合，且末端模之前的各项均为自由模，就称为这个末端模的\term[ziyou-xiaojie]{自由消解}[free resolution]。这里各 \(K_p=\bigwedge^pR^n\) 都是自由模，因此是否给出自由消解，关键在于尚待检验的正合性。

\ProseParagraphBreak
最短的例子是 \(n=1\)，复形就是 \(0\to R\xrightarrow{\times a_1}R\to0\)。若 \(k\) 为域，\(R=k[\varepsilon]/(\varepsilon^2)\)、\(a_1=\varepsilon\)，则左侧映射的核为非零理想 \((\varepsilon)\)，所以增广后也不是自由消解。哪些序列使增广后的 Koszul 复形成为 \(R/(a_1,\ldots,a_n)\) 的自由消解，将在第三卷第15.2—15.3节结合正则序列系统讨论。若一个生成元已经是单位，这里就能直接构造每个核元素的原像。

\begin{proposition}\label{b2:prop:koszul-unit-contractible}
设 \(R\) 为交换含幺环，\(n\ge1\)，\(a_1,\ldots,a_n\in R\)，且 \(a_1\) 是单位。令 \(e_1,\ldots,e_n\) 为 \(R^n\) 的标准基，\(K_\bullet=K_\bullet(a_1,\ldots,a_n;R)\) 为 Koszul 复形，其微分由 \(R\) 线性映射 \(\lambda:R^n\rightarrow R\)、\(\lambda(e_i)=a_i\) 的收缩给出。把 \(K_p=\bigwedge^pR^n\) 在 \(p<0\) 或 \(p>n\) 时取为零。定义次数为 \(+1\) 的 \(R\) 线性映射
\[
h_p:K_p\longrightarrow K_{p+1},
\qquad h_p(\xi)=a_1^{-1}e_1\wedge\xi.
\]
则在每个 \(0\le p\le n\) 上有
\[
\mathrm{d}_{p+1}h_p+h_{p-1}\mathrm{d}_p=\operatorname{id}_{K_p},
\]
其中边界处的映射取零。因此 \(K_\bullet\) 正合。
\end{proposition}
\begin{proof}
由\cref{b2:prop:exterior-contraction}的分次 Leibniz 规则及 \(\mathrm{d}_1(e_1)=a_1\)，对 \(\xi\in K_p\) 有
\[
\begin{aligned}
\mathrm{d}_{p+1}h_p(\xi)
&=a_1^{-1}\cdot\mathrm{d}_{p+1}(e_1\wedge\xi)\\
&=a_1^{-1}\bigl(a_1\xi-e_1\wedge\mathrm{d}_p(\xi)\bigr)\\
&=\xi-h_{p-1}\mathrm{d}_p(\xi).
\end{aligned}
\]
当 \(p=0\) 时，\(\mathrm{d}_0=0\)；当 \(p=n\) 时，\(e_1\wedge\xi\in\bigwedge^{n+1}R^n=0\)，上述等式在两个端点仍成立。这就证明所述恒等式，通常简写为 \(\mathrm{d}h+h\mathrm{d}=\operatorname{id}\)。

若 \(\mathrm{d}_p(\xi)=0\)，恒等式给出
\(\xi=\mathrm{d}_{p+1}(h_p(\xi))\)。这为每个闭元素（微分为零的元素）明确给出了成为边界（下一次微分的像）的原像。故 \(\ker\mathrm{d}_p\subseteq\operatorname{im}\mathrm{d}_{p+1}\)，反向包含由微分平方为零成立，复形在每处正合。
\end{proof}


\begin{chaptersummary}
交替性要求重复输入时取值为零，比单纯的交换变号更能适应一般系数环。外幂通过商掉重复因子的关系表达这一性质；有限自由模的递增指标楔积构成基，证明只用交替求和和系数检测，因而不要求阶乘可逆。

\ProseParagraphBreak
楔积把所有外幂连成分次代数，分次交换律记录因子换位的次数。直和的外代数分解须在张量乘法中保留这一符号。最高外幂只有一个基向量，使线性自同态表现为乘以行列式；较低次外幂则同时记录各阶子式，函子复合给出 Cauchy–Binet 公式。

\ProseParagraphBreak
伴随矩阵恒等式不仅给出可逆判据，也通过系数比较证明一般交换环上的 Cayley–Hamilton 定理。在域上，中间次数的非零纯楔积记录子空间；四维空间中的一条 Plücker 二次关系判定二向量是否具有这种形式，特征二时不能改用平方为零来判定。整个外代数上的收缩则把生成元的两两关系及其相容性组成 Koszul 复形。收缩平方为零保证相邻箭头复合为零，单位条件给出一个可直接证明正合的特例，一般序列的正合性仍需另行检查。
\end{chaptersummary}

\BookExercises

\begin{exercise}\label{b2:ex:09-alternating-characteristic-two}
在 \(\mathbb F_2^2\) 上分别写出全部反对称双线性形式和全部交替双线性形式的矩阵。两类形式各有多少个？
\end{exercise}
\begin{solution}
特征二使反对称条件成为 \(A^{\mathsf t}=A\)，所以矩阵为 \(\begin{pmatrix}a&b\\b&c\end{pmatrix}\)，三个参数独立，有 \(2^3=8\) 个。交替还要求对角线为零，所以矩阵为 \(\begin{pmatrix}0&b\\b&0\end{pmatrix}\)，只有两个。充分性可直接展开 \(x^{\mathsf t}Ax\)：交替情形下两个交叉项之和为零，对角项已经为零。\(A=I_2\) 是反对称而不交替的例子，因为在 \(e_1\) 上的自配对为一。
\end{solution}

\begin{exercise}\label{b2:ex:09-skew-not-exterior}
将 \(M=\mathbb Z/6\mathbb Z\)、\(L=\mathbb Z/4\mathbb Z\) 视为 \(\mathbb Z\) 模。求出全部双线性映射 \(M\times M\to L\)，判断哪些反对称、哪些交替，并说明哪些能够通过 \(\bigwedge^2M\) 分解。
\end{exercise}
\begin{solution}
双线性映射由 \(\ell=b(\bar1,\bar1)\in L\) 决定，公式为 \(b(\bar a,\bar c)=ac\ell\)。代表元改变六的倍数时，该式良定义恰要求 \(6\ell=0\)。在 \(\mathbb Z/4\mathbb Z\) 中，这等价于 \(\ell=\bar0\) 或 \(\bar2\)，所以共有两个映射。二者都满足 \(2\ell=0\)，从而 \(b(\bar c,\bar a)=-b(\bar a,\bar c)\)，均反对称；交替条件却要求 \(b(\bar1,\bar1)=0\)，故只有零映射交替。因为 \(M\) 由一个元素生成，\(\bigwedge^2M=0\)，也只有零映射能经过它分解。非零映射再次显示：目标中的二挠可以使交换变号成立，而重复输入仍不为零。
\end{solution}

\begin{exercise}\label{b2:ex:09-alternating-top-form}
设 \(R\) 为交换环，\(L\) 为 \(R\) 模，\(b:(R^3)^3\to L\) 为交替三线性映射，且 \(b(e_1,e_2,e_3)=\ell\)。证明 \(b\) 由 \(\ell\) 唯一决定，并计算 \(b(e_1+e_2,e_2+e_3,e_3+e_1)\)。
\end{exercise}
\begin{solution}
由\cref{b2:prop:exterior-alternating-universal}，\(b\) 对应线性映射 \(\bigwedge^3R^3\to L\)。该外幂以 \(e_1\wedge e_2\wedge e_3\) 为单个基向量，因此由其像 \(\ell\) 唯一决定。题给三向量楔积展开后，只有 \(e_1\wedge e_2\wedge e_3\) 与 \(e_2\wedge e_3\wedge e_1\) 两项没有重复因子；后三循环符号为正，所以结果为 \(2\ell\)。这个答案在有二挠时也成立，不能预先把它当作非零。
\end{solution}

\begin{exercise}\label{b2:ex:09-cyclic-exterior}
设 \(I\) 为交换环 \(R\) 的理想，\(M=R/I\)。计算 \(\bigwedge^pM\) 的全部次数，并证明
\[
\bigwedge M\cong R[x]/(x^2,Ix)
\]
为分次 \(R\) 代数同构，其中 \(x\) 次数为一。
\end{exercise}
\begin{solution}
\(M\) 由 \(\bar1\) 生成，任意至少二重的纯楔积都是 \(\bar1\wedge\bar1\wedge\cdots\) 的标量倍数，所以为零。因此零次分量为 \(R\)，一次分量为 \(R/I\)，更高分量为零。右侧商代数的元素唯一写成常数 \(r\) 加一次项 \(\bar a x\)，乘法由 \(x^2=0\) 与 \(Ix=0\) 决定。将 \(\bar1\in M\) 送到 \(x\) 给出双向保持这些关系的映射，或逐次数比较上述唯一表达，即得同构。
\end{solution}

\begin{exercise}\label{b2:ex:09-wedge-signs}
在 \(\bigwedge\mathbb Z^4\) 中计算
\[
(e_1+e_3)\wedge(e_2-e_4),\qquad
(e_1\wedge e_3)\wedge(e_2\wedge e_4),\qquad
(e_1\wedge e_2+e_3\wedge e_4)^2.
\]
说明为什么“正次数元素都平方为零”是错误的。
\end{exercise}
\begin{solution}
第一项为 \(e_1\wedge e_2-e_1\wedge e_4-e_2\wedge e_3-e_3\wedge e_4\)。第二项把 \(e_3\) 与 \(e_2\) 交换一次，得到 \(-e_1\wedge e_2\wedge e_3\wedge e_4\)。第三项的两个同侧平方为零，两个交叉项因次数均为二而同号，故结果为 \(2e_1\wedge e_2\wedge e_3\wedge e_4\)，在整数系数下非零。\cref{b2:prop:exterior-graded-commutativity}保证奇数次齐次元素平方为零，并不对所有偶数次元素作同样断言。
\end{solution}

\begin{exercise}\label{b2:ex:09-wedge-independence}
设 \(A:k^n\to k^m\) 为域上的线性映射，\(p\ge1\)。证明 \(\bigwedge^pA\ne0\) 当且仅当 \(\operatorname{rank}A\ge p\)，并用 \(p\) 阶子式表述这个条件。对于
\[
A=\begin{pmatrix}1&0&1\\0&1&1\\1&1&2\end{pmatrix},
\]
在任意域上求最大的 \(p\)，使 \(\bigwedge^pA\ne0\)。
\end{exercise}
\begin{solution}
若 \(\operatorname{rank}A<p\)，任意 \(p\) 个像向量都相关，\cref{b2:prop:wedge-linear-independence}说明每个纯楔积都被 \(\bigwedge^pA\) 送到零；纯楔积生成整个外幂，所以该映射为零。若秩至少为 \(p\)，可选 \(p\) 个向量，使它们的像独立，其楔积非零，故 \(\bigwedge^pA\ne0\)。当 \(p\le\min(m,n)\) 时，由\eqref{b2:eq:exterior-minor-action}，该条件等价于至少一个 \(p\) 阶子式非零；超过这两个维数之一时，外幂映射为零，也没有这样的子式。

题给矩阵的第三列是前两列之和，而左上二阶子式为 \(1\)。因此它在任意域上秩均为二，最大的 \(p\) 是二。这里既用一个非零子式给出秩的下界，也用列关系给出上界，不能只由某个二阶子式非零就断言秩恰为二。
\end{solution}

\begin{exercise}\label{b2:ex:09-exterior-square-map}
将整数矩阵 \(A=\begin{pmatrix}1&1&0\\0&2&1\\0&0&3\end{pmatrix}\) 视为 \(\mathbb Z^3\) 上的线性映射。在基 \(e_1\wedge e_2,e_1\wedge e_3,e_2\wedge e_3\) 下计算 \(\bigwedge^2 A\) 的矩阵，并比较其行列式与 \((\det A)^2\)。
\end{exercise}
\begin{solution}
三列的像分别为
\(2e_1\wedge e_2\)、\(e_1\wedge e_2+3e_1\wedge e_3\)、\(e_1\wedge e_2+3e_1\wedge e_3+6e_2\wedge e_3\)。所以矩阵为
\[
\begin{pmatrix}2&1&1\\0&3&3\\0&0&6\end{pmatrix}.
\]
其行列式为 \(36\)，而 \(\det A=6\)，在本例中确有 \(\det(\bigwedge^2A)=(\det A)^2\)。这些数值是整数恒等式，向任意交换环映射后仍然成立；这里只核验本例，不以它代替一般公式的证明。
\end{solution}

\begin{exercise}\label{b2:ex:09-odd-square}
设 \(R\) 为交换环，\(M\) 为具有 \(n\) 个基向量的自由 \(R\) 模，\(N=\bigoplus_{p\ge1}\bigwedge^pM\) 为外代数中的正次理想。证明 \(N^{n+1}=0\)，并证明元素 \(a_0+\eta\)、\(a_0\in R,\eta\in N\)，可逆当且仅当 \(a_0\) 是 \(R\) 的单位。写出逆元的有限求和公式。
\end{exercise}
\begin{solution}
每个正次齐次分量的次数至少为一，所以 \(n+1\) 个正次元素相乘的所有分量都在次数 \(n+1\) 以上，由\cref{b2:thm:exterior-power-basis}全部为零。因此 \(N^{n+1}=0\)。投影到零次分量是保持单位元的代数同态 \(\bigwedge M\to R\)，故可逆元素的常数项也必须可逆。

反过来，若 \(a_0\) 可逆，令 \(z=a_0^{-1}\eta\)。此时 \(z^{n+1}=0\)，有限几何级数给出
\[
(a_0+\eta)^{-1}=a_0^{-1}\sum_{j=0}^n(-z)^j.
\]
左乘或右乘 \(a_0+\eta=a_0(1+z)\)，相邻项均抵消，只剩 \(1+(-1)^n z^{n+1}=1\)。这里每项都是同一个元素 \(z\) 的幂，且标量 \(a_0\) 在中心，所以即使外代数并非普通交换代数，仍得到双侧逆。
\end{solution}

\begin{exercise}\label{b2:ex:09-direct-sum-degree-three}
设 \(R\) 为交换环，\(M,N\) 为 \(R\) 模，分别有基 \(e_1,e_2\) 与 \(f_1,f_2\)。用直和外幂分解写出 \(\bigwedge^3(M\oplus N)\) 的一组基，并说明 \((0,f_1)\wedge(e_1,0)\wedge(0,f_2)\) 对应哪个张量及其符号。
\end{exercise}
\begin{solution}
因为两侧的三次外幂均为零，只剩
\[
\left(\bigwedge^2M\otimes_RN\right)
\oplus\left(M\otimes_R\bigwedge^2N\right).
\]
一组基为 \((e_1\wedge e_2)\otimes f_1\)、\((e_1\wedge e_2)\otimes f_2\)、\(e_1\otimes(f_1\wedge f_2)\)、\(e_2\otimes(f_1\wedge f_2)\)。题给楔积要把 \(e_1\) 移到最前，交换一次，故对应 \(-e_1\otimes(f_1\wedge f_2)\)。若系数环特征二，负号等于正号，但符号产生的规则不变。
\end{solution}

\begin{exercise}\label{b2:ex:09-exterior-not-injective}
设 \(k\) 为域，\(R=k[x,y]\)、\(I=(x,y)\)。进一步计算正文反例中的 \(\bigwedge^2I\)，证明
\[
R/I\longrightarrow\bigwedge^2I,\qquad \bar r\longmapsto r(x\wedge y)
\]
是 \(R\) 模同构，并求包含 \(I\hookrightarrow R\) 所诱导的二次外幂映射的核。
\end{exercise}
\begin{solution}
因为 \(I\) 由 \(x,y\) 生成，双线性展开说明 \(\bigwedge^2I\) 由 \(w=x\wedge y\) 生成。又
\[
xw=x\wedge(xy)=x\wedge(yx)=y(x\wedge x)=0,
\qquad yw=(xy)\wedge y=x(y\wedge y)=0.
\]
所以 \(Iw=0\)，题给映射良定义且满射。\cref{b2:prop:exterior-injection-failure}中由一次项构造的交替映射诱导 \(\widetilde\beta:\bigwedge^2I\to k=R/I\)，并满足 \(\widetilde\beta(w)=1\)。它与题给映射的复合是 \(R/I\) 的恒等，故题给映射也单射。由于 \(\bigwedge^2R=0\)，包含所诱导的映射是零映射，其核为整个 \(\bigwedge^2I\cong R/I\)。这个例子不仅检测到一个非零核元，还算出了全部核。
\end{solution}

\begin{exercise}\label{b2:ex:09-determinant-basis-change}
设 \(A=(a_{ij})\) 为交换环上的 \(n\times n\) 上三角矩阵，\(1\le p\le n\)。把递增指标组 \(I=(i_1,\ldots,i_p)\) 按字典序排列。证明 \(\bigwedge^pA\) 在相应外幂基下仍为上三角矩阵，求出其对角元，并证明
\[
\det\left(\bigwedge^pA\right)=(\det A)^{\binom{n-1}{p-1}}.
\]
\end{exercise}
\begin{solution}
\(Ae_j\) 只涉及 \(e_1,\ldots,e_j\)。若 \(e_I\) 在 \(Ae_{j_1}\wedge\cdots\wedge Ae_{j_p}\) 中出现非零系数，则必有 \(i_s\le j_s\) 对所有 \(s\) 成立：否则前 \(s\) 个输入所选的不同基指标都须不超过 \(j_s\)，但目标指标组中这样的指标至多有 \(s-1\) 个，不可能。于是非零系数只能出现于字典序不大于 \(J\) 的行 \(I\)，证明矩阵上三角。

对应 \(I=J\) 的对角元是子矩阵 \(A_{J,J}\) 的行列式，而它仍上三角，故为 \(\prod_{j\in J}a_{jj}\)。于是
\[
\det\left(\bigwedge^pA\right)
=\prod_{|J|=p}\prod_{j\in J}a_{jj}
=\prod_{j=1}^n a_{jj}^{\binom{n-1}{p-1}}
=(\det A)^{\binom{n-1}{p-1}}.
\]
中间的指数来自固定 \(j\) 后，从其余 \(n-1\) 个指标中选取 \(p-1\) 个的方式数；没有使用可逆性或约去任何对角元。
\end{solution}

\begin{exercise}\label{b2:ex:09-gram-cauchy-binet}
设 \(A\) 为实 \(2\times n\) 矩阵。证明
\[
\det(AA^{\mathsf t})=\sum_{1\le i<j\le n}\det\bigl(A_{\{1,2\},\{i,j\}}\bigr)^2.
\]
由此证明该行列式非负，并说明它何时为正。
\end{exercise}
\begin{solution}
将\cref{b2:thm:cauchy-binet}用于 \(A\) 与 \(A^{\mathsf t}\)。每个对应的行子式是原列子式的转置，行列式相同，因此每项是平方。和非负，且为正当且仅当至少一个二阶子式非零。若有两个独立列，对应二阶子式非零；若没有，全部列至多生成一维空间。故为正恰好等价于 \(A\) 的秩为二。\(n<2\) 时为空和，行列式为零，结论仍成立。
\end{solution}

\begin{exercise}\label{b2:ex:09-block-triangular}
对交换环上的分块矩阵 \(C=\begin{pmatrix}A&B\\0&D\end{pmatrix}\)，其中 \(A,D\) 分别为 \(m,n\) 阶方阵，证明 \(\det C=\det A\det D\)，不要求两个块可逆。
\end{exercise}
\begin{solution}
将 \(R^{m+n}\) 的基分为前 \(m\) 个 \(e_i\) 和后 \(n\) 个 \(f_j\)。前 \(m\) 列的楔积为 \((\det A)e_1\wedge\cdots\wedge e_m\)。在其后乘上后 \(n\) 列的楔积时，只要从任何一列选到前 \(m\) 个坐标的分量，就与前面的全部 \(e_i\) 重复而为零。因此只有后块 \(D\) 的分量贡献，其系数为 \(\det D\)，得到结论。这个推理只用展开和重复因子为零，不需要消去 \(\det A\) 或取逆。
\end{solution}

\begin{exercise}\label{b2:ex:09-adjugate-mod-eight}
在 \(R=\mathbb Z/8\mathbb Z\) 上判断 \(A=\begin{pmatrix}1&2\\3&1\end{pmatrix}\) 是否可逆；若可逆，用伴随矩阵求逆。再判断 \(B=\begin{pmatrix}2&0\\0&1\end{pmatrix}\) 是否可逆。
\end{exercise}
\begin{solution}
\(\det A=1-6=\bar3\)，它的逆为 \(\bar3\)，因此
\[
A^{-1}=3\begin{pmatrix}1&-2\\-3&1\end{pmatrix}
=\begin{pmatrix}3&2\\7&3\end{pmatrix}\pmod8.
\]
直接相乘也得到单位矩阵。\(\det B=\bar2\) 非零但不是单位，所以 \(B\) 不可逆；事实上 \(B(\bar4,\bar0)=(\bar0,\bar0)\)，连单射都不满足。
\end{solution}

\begin{exercise}\label{b2:ex:09-cayley-recurrence}
令 \(A=\begin{pmatrix}1&1\\1&0\end{pmatrix}\)，\(F_0=0,F_1=1,F_{n+1}=F_n+F_{n-1}\)。用 Cayley–Hamilton 定理证明 \(A^n=F_nA+F_{n-1}I_2\) 对所有 \(n\ge1\) 成立，并求 \(A^5\)。
\end{exercise}
\begin{solution}
\(\chi_A(t)=t^2-t-1\)，故 \(A^2=A+I_2\)。\(n=1\) 的公式成立；若 \(n\) 时成立，则
\(A^{n+1}=F_n(A+I_2)+F_{n-1}A=F_{n+1}A+F_nI_2\)，归纳完成。\(F_5=5,F_4=3\)，所以 \(A^5=5A+3I_2=\begin{pmatrix}8&5\\5&3\end{pmatrix}\)。这是整数矩阵恒等式，在任意交换环中解释相应整数系数后仍成立。
\end{solution}

\begin{exercise}\label{b2:ex:09-characteristic-coefficients}
设 \(A\in M_n(R)\)、\(n\ge1\)，其中 \(R\) 交换。证明 \(\chi_A(t)\) 的 \(t^{n-1}\) 系数为 \(-\operatorname{tr}A\)，常数项为 \((-1)^n\det A\)。说明该论证为什么不需要特征值。
\end{exercise}
\begin{solution}
常数项为 \(\det(-A)=(-1)^n\det A\)，由各列提出负号得到。为得到 \(t^{n-1}\)，Leibniz 展开中须选取 \(n-1\) 个对角线上的 \(t\)；相应置换已固定 \(n-1\) 个指标，因而也固定最后一个，所以只能来自 \(\prod_i(t-a_{ii})\)。从第 \(i\) 项选 \(-a_{ii}\)、其余选 \(t\)，求和得到 \(-\sum_i a_{ii}\)。全程是交换环上的多项式系数计算，不要求特征多项式分裂或存在特征值。
\end{solution}

\begin{exercise}\label{b2:ex:09-plucker-explicit}
在任意域上，判断坐标按 \(12,13,14,23,24,34\) 排列为 \((1,2,0,3,1,2)\) 的二张量是否可分解；若是，给出两个因子。这里整数按基域的特征解释。
\end{exercise}
\begin{solution}
关系左端为 \(1\cdot2-2\cdot1+0\cdot3=0\)，所以\cref{b2:prop:plucker-four-dimensional}适用。可取
\[
u=e_1-3e_3-e_4,\qquad v=e_2+2e_3.
\]
展开后五个前项坐标分别为 \(1,2,0,3,1\)，\(34\) 坐标为 \((-3)\cdot0-(-1)\cdot2=2\)。因为 \(12\) 坐标始终为一，楔积在任何特征下都非零，两个因子均线性无关。
\end{solution}

\begin{exercise}\label{b2:ex:09-plucker-plane-recovery}
设 \(k\) 为域，\(e_1,\ldots,e_4\) 为 \(k^4\) 的标准基。令 \(\omega=(e_1+e_3)\wedge(e_2+e_4)\in\bigwedge^2k^4\)。求所有满足 \(v\wedge\omega=0\) 的向量 \(v\)，并说明它与 \(\omega\) 的两条生成向量之间的关系。
\end{exercise}
\begin{solution}
写 \(v=\sum_i a_ie_i\)。因为 \(e_1+e_3,e_2+e_4\) 独立，\eqref{b2:eq:plucker-recovers-plane}说明解空间正是它们的线性生成空间，所以
\[
v=a(e_1+e_3)+b(e_2+e_4),\qquad a,b\in k.
\]
等价条件为 \(a_3=a_1\)、\(a_4=a_2\)。也可把 \(\omega=e_1\wedge e_2+e_1\wedge e_4-e_2\wedge e_3+e_3\wedge e_4\) 代入，再比较四个递增三重楔积的系数，得到同样的两个独立方程。
\end{solution}

\begin{exercise}\label{b2:ex:09-contractions-anticommute}
设 \(R\) 为交换环，\(M\) 为 \(R\) 模，\(m\in M\)、\(\lambda:M\to R\) 线性。令 \(L_m(\xi)=m\wedge\xi\)。证明
\[
\iota_\lambda L_m+L_m\iota_\lambda=\lambda(m)\operatorname{id}_{\bigwedge M}.
\]
当 \(M=R^2\)、\(m=e_1\)、\(\lambda(e_1)=1\)、\(\lambda(e_2)=0\) 时，在基 \(1,e_1,e_2,e_1\wedge e_2\) 下写出两个算子的矩阵并核对等式。
\end{exercise}
\begin{solution}
收缩的分次乘法公式给出
\[
\iota_\lambda(m\wedge\xi)=\lambda(m)\xi-m\wedge\iota_\lambda(\xi),
\]
移项即得等式，无须除以二。在指定基下，两个矩阵分别为
\[
[L_{e_1}]=\begin{pmatrix}0&0&0&0\\1&0&0&0\\0&0&0&0\\0&0&1&0\end{pmatrix},\qquad
[\iota_\lambda]=\begin{pmatrix}0&1&0&0\\0&0&0&0\\0&0&0&1\\0&0&0&0\end{pmatrix}.
\]
前者将 \(1,e_2\) 分别送到 \(e_1,e_1\wedge e_2\)，后者将 \(e_1,e_1\wedge e_2\) 分别送到 \(1,e_2\)，其余基向量都送到零。因此两种复合的和在四个基向量上均为恒等。这条混合关系与\cref{b2:cor:contractions-anticommute}中的两个收缩反交换关系不同：它的右端可以非零。
\end{solution}

\begin{exercise}\label{b2:ex:09-koszul-unit}
设 \(a_1,\ldots,a_n\) 生成交换环 \(R\) 的单位理想，取 \(b_i\in R\) 使 \(\sum_i b_i a_i=1\)。在 Koszul 复形中令 \(u=\sum_i b_ie_i\)、\(h(\xi)=u\wedge\xi\)，证明 \(\mathrm{d}h+h\mathrm{d}=\operatorname{id}\)。对 \(R=\mathbb Z\)、序列 \(2,3\)，显式写出各次 \(h\) 的矩阵，说明正合性不要求序列中一定含有单位。
\end{exercise}
\begin{solution}
因为 \(\mathrm{d}(u)=\sum_i b_i a_i=1\)，收缩的分次乘法公式给出 \(\mathrm{d}(u\wedge\xi)=\xi-u\wedge\mathrm{d}\xi\)，故所求等式成立。若 \(\mathrm{d}\xi=0\)，便有 \(\xi=\mathrm{d}(h\xi)\)；配合 \(\mathrm{d}^2=0\)，得到每个位置正合。

在整数例中，取 \(u=-e_1+e_2\)，因为 \(-2+3=1\)。在 \(K_1\) 的基 \(e_1,e_2\) 和 \(K_2\) 的基 \(e_1\wedge e_2\) 下，
\[
\mathrm{d}_1=\begin{pmatrix}2&3\end{pmatrix},\quad
\mathrm{d}_2=\binom{-3}{2},\quad
h_0=\binom{-1}{1},\quad h_1=\begin{pmatrix}-1&-1\end{pmatrix},\quad h_2=0.
\]
直接相乘得 \(\mathrm{d}_1h_0=1\)、\(\mathrm{d}_2h_1+h_0\mathrm{d}_1=I_2\)、\(h_1\mathrm{d}_2=1\)。所以该复形正合，尽管 \(2,3\) 在 \(\mathbb Z\) 中都不是单位。\cref{b2:prop:koszul-unit-contractible}的单位元条件是充分条件；本题用 Bézout 组合将它推广到生成单位理想的序列。
\end{solution}

\begin{exercise}\label{b2:ex:09-koszul-three}
设 \(k\) 为任意域，令 \(R=k[t]\)，考虑序列 \(t,t,t\) 的 Koszul 复形。在 \(K_1\) 的基 \(e_1,e_2,e_3\) 与 \(K_2\) 的基 \(e_1\wedge e_2,e_1\wedge e_3,e_2\wedge e_3\) 下写出全部微分矩阵，核对相邻乘积为零，并计算
\[
H_p\coloneqq\ker\mathrm{d}_p/\operatorname{im}\mathrm{d}_{p+1}\qquad(0\le p\le3),
\]
其中 \(\mathrm{d}_0=0\)、\(K_4=0\)。这些商模称为复形的同调模；它们如何判断各位置的正合性？
\end{exercise}
\begin{solution}
微分矩阵为
\[
\mathrm{d}_1=t\begin{pmatrix}1&1&1\end{pmatrix},\qquad
\mathrm{d}_2=t\begin{pmatrix}-1&-1&0\\1&0&-1\\0&1&1\end{pmatrix},\qquad
\mathrm{d}_3=t\begin{pmatrix}1\\-1\\1\end{pmatrix}.
\]
两个相邻乘积都为零。由于 \(R\) 是整环，求核时可消去非零的 \(t\)。令 \(u=(-1,1,0)\)、\(v=(-1,0,1)\)，则
\[
\ker\mathrm{d}_1=Ru\oplus Rv,\qquad
\operatorname{im}\mathrm{d}_2=tRu\oplus tRv.
\]
后一等式来自中间矩阵的三列 \(u,v,v-u\)。同样，\(\ker\mathrm{d}_2=R(1,-1,1)\)，而 \(\operatorname{im}\mathrm{d}_3=tR(1,-1,1)\)；\(\mathrm{d}_3\) 单射。因此
\[
H_0\cong R/(t),\qquad H_1\cong(R/(t))^2,\qquad
H_2\cong R/(t),\qquad H_3=0.
\]
在第 \(p\) 项正合恰好等价于 \(H_p=0\)。这个复形在 \(K_1,K_2\) 处不正合，尽管两个相邻乘积都为零；接上末端 \(R/(t)\) 只能修复零次余核，不能消去这两个非零同调模。
\end{solution}
